% ===== preamble =====
\usepackage[T1]{fontenc}\usepackage{lmodern}
\usepackage[margin=0.9in]{geometry}
\usepackage{amsmath,amssymb,amsthm,mathtools}
\usepackage{mathrsfs}
\usepackage{enumitem}
\usepackage{microtype}
\usepackage{tikz}
\usetikzlibrary{automata,positioning,arrows.meta,calc,fit,chains,
  shapes.geometric,shapes.misc,decorations.pathreplacing,decorations.markings,
  backgrounds,matrix,patterns}
\usepackage{tikz-cd}
\usepackage[hidelinks]{hyperref}

\setlength{\parindent}{0pt}\setlength{\parskip}{3pt}
% absorb hard-to-break lines (long inline-math lists) without global \sloppy
\emergencystretch=3em
\makeatletter
\renewcommand*\l@section{\@dottedtocline{1}{0pt}{2.8em}}
% running head: use the (short) chapter mark, do not let it overflow the line
\renewcommand{\chaptermark}[1]{\markboth{\MakeUppercase{\chaptername\ \thechapter.\ #1}}{}}
\makeatother

% ---- theorem environments (results numbered chapter.section.k) ----
\theoremstyle{plain}
\newtheorem{theorem}{Theorem}[section]
\newtheorem{proposition}[theorem]{Proposition}
\newtheorem{lemma}[theorem]{Lemma}
\newtheorem{corollary}[theorem]{Corollary}
\theoremstyle{definition}
\newtheorem{definition}[theorem]{Definition}
\newtheorem{example}[theorem]{Example}
\theoremstyle{remark}
\newtheorem*{remark}{Remark}

% ---- worked example (clean, no box; shares the result counter) ----
\newenvironment{wex}[1]{%
  \par\addvspace{\medskipamount}\refstepcounter{theorem}%
  \noindent\textbf{Worked Example~\thetheorem.}\enspace\textit{#1.}\par\nobreak\smallskip}%
  {\par\addvspace{\medskipamount}}
\newcommand{\wrec}{\par\smallskip\noindent\textit{Recognition.}\enspace\ignorespaces}
\newcommand{\wsol}{\par\smallskip\noindent\textit{Solution.}\enspace\ignorespaces}
\newcommand{\whichrule}[1]{\par\smallskip\noindent$\blacktriangleright$\ \textbf{Which idea, and why?}\enspace #1\par}

% ---- method / strategy block (rule-delimited, no color) ----
\newenvironment{method}[1]{%
  \par\addvspace{\medskipamount}\noindent\rule{\linewidth}{0.4pt}\par\nobreak\vskip2pt
  \noindent\textbf{Method.}\enspace\textit{#1}\par\nobreak\smallskip}%
  {\par\nobreak\vskip2pt\noindent\rule{\linewidth}{0.4pt}\par\addvspace{\medskipamount}}

% ---- caution note (run-in, no color) ----
\newenvironment{pitfall}{%
  \par\addvspace{\smallskipamount}\noindent\textbf{\textit{Caution.}}\enspace\ignorespaces}%
  {\par\addvspace{\smallskipamount}}

% ---- exercise list ----
\newenvironment{exers}[1]{\par\medskip\noindent\textbf{#1}\par\nopagebreak
  \begin{enumerate}[leftmargin=*,labelsep=0.6em,itemsep=2pt,topsep=2pt]}
  {\end{enumerate}}

% ================= math shortcuts: Theory of Computation =================
% ---- number systems / common sets ----
\newcommand{\R}{\mathbb{R}}\newcommand{\Q}{\mathbb{Q}}\newcommand{\Z}{\mathbb{Z}}
\newcommand{\N}{\mathbb{N}}\newcommand{\C}{\mathbb{C}}\newcommand{\F}{\mathbb{F}}
\newcommand{\es}{\varnothing}
\newcommand{\set}[1]{\left\{#1\right\}}
\newcommand{\abs}[1]{\left|#1\right|}
\newcommand{\card}[1]{\left|#1\right|}
\newcommand{\st}{\mid}
\newcommand{\imp}{\implies}\newcommand{\eqv}{\iff}
\newcommand{\inv}{^{-1}}
\providecommand{\id}{\operatorname{id}}

% ---- alphabets, strings, languages ----
\newcommand{\Sig}{\Sigma}
\newcommand{\Gam}{\Gamma}
\newcommand{\eps}{\varepsilon}
\newcommand{\str}[1]{\mathtt{#1}}
\newcommand{\lang}[1]{L(#1)}
\newcommand{\Lang}{\mathcal{L}}
\newcommand{\Sigstar}{\Sigma^{*}}
\newcommand{\Gamstar}{\Gamma^{*}}
\newcommand{\rev}[1]{#1^{\mathcal{R}}}
\newcommand{\concat}{\cdot}
\newcommand{\kstar}{^{*}}
\newcommand{\kplus}{^{+}}
\newcommand{\substr}{\sqsubseteq}

% ---- automata data ----
\newcommand{\trans}{\delta}
\newcommand{\xtrans}{\widehat{\delta}}
\newcommand{\qstart}{q_0}\newcommand{\Acc}{F}
\newcommand{\pow}[1]{\mathcal{P}(#1)}
\newcommand{\dfa}{M}\newcommand{\nfa}{N}
\newcommand{\yields}{\vdash}
\newcommand{\yieldstar}{\vdash^{*}}

% ---- regular expressions ----
\newcommand{\reor}{\cup}
\newcommand{\restar}{^{*}}

% ---- language / complexity classes ----
\newcommand{\REG}{\mathsf{REGULAR}}
\newcommand{\CFL}{\mathsf{CFL}}
\newcommand{\classP}{\mathsf{P}}
\newcommand{\classNP}{\mathsf{NP}}
\newcommand{\coNP}{\mathsf{coNP}}
\newcommand{\PSPACE}{\mathsf{PSPACE}}
\newcommand{\NPSPACE}{\mathsf{NPSPACE}}
\newcommand{\EXP}{\mathsf{EXPTIME}}
\newcommand{\Lclass}{\mathsf{L}}
\newcommand{\NL}{\mathsf{NL}}
\newcommand{\TIME}{\mathrm{TIME}}\newcommand{\SPACE}{\mathrm{SPACE}}
\newcommand{\NTIME}{\mathrm{NTIME}}\newcommand{\NSPACE}{\mathrm{NSPACE}}

% ---- decision problems (Sipser bracket notation) ----
\newcommand{\enc}[1]{\langle #1\rangle}
\newcommand{\ATM}{A_{\mathsf{TM}}}
\newcommand{\HALT}{\mathit{HALT}_{\mathsf{TM}}}
\newcommand{\ETM}{E_{\mathsf{TM}}}
\newcommand{\ADFA}{A_{\mathsf{DFA}}}\newcommand{\ANFA}{A_{\mathsf{NFA}}}
\newcommand{\ACFG}{A_{\mathsf{CFG}}}
\newcommand{\SAT}{\mathit{SAT}}\newcommand{\CSAT}{\mathit{3SAT}}
\newcommand{\CLIQUE}{\mathit{CLIQUE}}\newcommand{\HAMPATH}{\mathit{HAMPATH}}
\newcommand{\VC}{\mathit{VERTEX\text{-}COVER}}
% ---- TM tape blank + emptiness/equivalence problems + decidability classes ----
% (union of the per-chapter preambles: \blank from Ch.8; \E,\EQTM,\REC,\DEC from Ch.9)
\newcommand{\blank}{\sqcup}                        % TM tape blank (Sipser open box)
\newcommand{\E}{E}                                % emptiness-problem letter, e.g. \E_{\mathsf{DFA}}
\newcommand{\EQTM}{\mathit{EQ}_{\mathsf{TM}}}
\newcommand{\REC}{\mathsf{REC}}                   % Turing-recognizable class
\newcommand{\DEC}{\mathsf{DEC}}                   % decidable class

% ---- reductions / asymptotics ----
\newcommand{\redm}{\leq_{\mathrm{m}}}
\newcommand{\redp}{\leq_{\mathrm{p}}}
\newcommand{\redT}{\leq_{\mathrm{T}}}
\newcommand{\redL}{\leq_{\mathrm{L}}}
\DeclareMathOperator{\bigO}{O}
\newcommand{\bigOmega}{\Omega}\newcommand{\bigTheta}{\Theta}
\DeclareMathOperator{\poly}{poly}

% ---- graph theory (for the foundations + reductions) ----
\newcommand{\Gr}{G}\newcommand{\Vset}{V}\newcommand{\Eset}{E}

% ================= DIAGRAM HOUSE STYLES =================
\tikzset{
  >={Stealth[length=2.2mm]},
  node distance=18mm and 22mm,
  fa/.style={->,shorten >=0.5pt,auto,semithick},
  every state/.style={draw=black,line width=0.7pt,minimum size=8.5mm,
    inner sep=1pt,font=\small},
  acc/.style={accepting by double},
  st/.style={state},
  stinit/.style={state,initial,initial text={}},
  stacc/.style={state,accepting by double},
  stinitacc/.style={state,initial,initial text={},accepting by double},
  tredge/.style={font=\small,inner sep=1.5pt},
  loop above/.append style={looseness=5},
  tape/.style={matrix of nodes,nodes={draw,minimum size=6.5mm,font=\small,
    anchor=center,inner sep=1pt},column sep=-\pgflinewidth,row sep=0pt},
  head/.style={-{Stealth[length=2mm]},line width=0.8pt},
  headlbl/.style={font=\footnotesize},
  treenode/.style={font=\small,inner sep=1.5pt},
  tedge/.style={line width=0.5pt,black!60},
  procbox/.style={draw,rounded corners,minimum height=8mm,minimum width=16mm,
    font=\small,inner sep=3pt,line width=0.6pt},
  flow/.style={-{Stealth[length=2.2mm]},line width=0.7pt},
  flowlbl/.style={font=\footnotesize,fill=white,inner sep=1.5pt},
  gv/.style={circle,fill=black,inner sep=1.7pt},
  gvlab/.style={circle,draw=black,line width=0.6pt,minimum size=6mm,
    inner sep=0.5pt,font=\small},
  gedge/.style={line width=0.7pt,black!75},
  gvlbl/.style={font=\footnotesize,fill=white,inner sep=1.4pt},
  ltdot/.style={circle,fill=black,inner sep=1.5pt},
  ltedge/.style={black!50,line width=0.7pt},
  ltlab/.style={font=\footnotesize,fill=white,inner sep=1.6pt},
}

% ================= additions to fold in the Integrated Framework appendix =================
\usepackage{adjustbox}
\usepackage{pdflscape}
\definecolor{cFound}{HTML}{4A4A4A}
\definecolor{cComb}{HTML}{1F6F8B}
\definecolor{cGraph}{HTML}{2E7D32}
\definecolor{cComp}{HTML}{8E44AD}
\definecolor{cBridge}{HTML}{B5651D}
\definecolor{cCore}{HTML}{8B0000}
\tikzset{
  clusterbox/.style={rounded corners=6pt,draw,line width=1pt,inner sep=10pt,
    fill=#1!4,draw=#1!70},
  clusterlbl/.style={font=\bfseries\small,text=#1!75},
  topic/.style={rounded corners=3pt,draw=#1!75,line width=0.8pt,fill=#1!6,
    text=black,align=center,font=\footnotesize,inner sep=3.4pt,
    minimum height=7mm,text width=26mm},
  topicw/.style={topic=#1,text width=34mm},
  core/.style={rounded corners=4pt,draw=cCore,line width=1.4pt,fill=cCore!8,
    text=black,align=center,font=\small\bfseries,inner sep=5pt,text width=52mm},
  egchip/.style={draw=black!35,line width=0.4pt,fill=black!3,rounded corners=2pt,
    font=\scriptsize,inner sep=2pt,align=center,text=black!75},
  pre/.style={-{Stealth[length=2.2mm]},line width=0.8pt,draw=black!70},
  xlink/.style={-{Stealth[length=2mm]},line width=0.6pt,draw=black!45,densely dashed},
  brid/.style={{Stealth[length=2.2mm]}-{Stealth[length=2.2mm]},line width=1pt,draw=cBridge!85},
  edgelbl/.style={font=\scriptsize,fill=white,inner sep=1.3pt,text=black!70},
}
\theoremstyle{definition}
\newtheorem{bridgeprob}{Bridge Problem}
\newtheorem{capstone}{Capstone}
\newcommand{\lkey}[2]{\tikz[baseline=-0.5ex]{\draw[#1] (0,0)--(0.7,0);}~\,{\footnotesize #2}}

% ---- structure-preserving squares, local concept maps, reduction graphs ----
% ---- commutative diagrams (tikz-cd line/tip styles; arrow LABELS given inline) ----
% usage: \arrow[r, pres, "{preserves }L"]  /  \arrow[d, quot]  /  \arrow[r, inc]
\tikzcdset{
  pres/.style={dashed},            % meaning/language-preserving transformation
  quot/.style={two heads},         % quotient / minimization (surjection)
  inc/.style={hook},               % inclusion / embedding (injection)
}

\tikzset{
  % ---- local concept-map nodes / edges (light grayscale) ----
  cnode/.style={rounded corners=2pt,draw=black!72,line width=0.6pt,align=center,
    font=\footnotesize,inner sep=3pt,text width=23mm,minimum height=7mm},
  cnodew/.style={cnode,text width=32mm},
  croot/.style={cnode,line width=1pt,draw=black,font=\footnotesize\bfseries,fill=black!4},
  cedge/.style={-{Stealth[length=2.2mm]},line width=0.6pt,draw=black!70},
  cbidir/.style={{Stealth[length=2.2mm]}-{Stealth[length=2.2mm]},line width=0.6pt,draw=black!70},
  cdash/.style={-{Stealth[length=2mm]},line width=0.6pt,draw=black!55,densely dashed},
  cedgelbl/.style={font=\scriptsize,fill=white,inner sep=1.2pt,text=black!75},
  % ---- structure-preserving square in plain TikZ (when tikz-cd is awkward) ----
  smnode/.style={align=center,font=\small,inner sep=2pt},
  smmap/.style={-{Stealth[length=2.4mm]},line width=0.7pt},
  smpres/.style={-{Stealth[length=2.4mm]},line width=0.9pt,densely dashed},
  smlbl/.style={font=\footnotesize,fill=white,inner sep=1.3pt},
  % ---- reduction-graph arrows: point FROM the known-hard problem TO the target ----
  redarrow/.style={-{Stealth[length=2.8mm]},line width=1pt,draw=black!82},
  redlbl/.style={font=\footnotesize,fill=white,inner sep=1.4pt},
  % ---- quotient / two-headed map in plain TikZ ----
  quotarrow/.style={-{Stealth[length=2.4mm]Stealth[length=2.4mm]},line width=0.8pt},
}

